\documentclass[11pt]{article}

\usepackage[margin=1in]{geometry}
\usepackage{amsmath, amssymb}
\usepackage{booktabs}
\usepackage{siunitx}
\usepackage[colorlinks=true, linkcolor=blue, urlcolor=blue]{hyperref}

\newcommand{\dpar}[2]{\frac{\partial #1}{\partial #2}}
\newcommand{\D}{\mathrm{d}}
\newcommand{\uvec}{\mathbf{u}}
\newcommand{\code}[1]{\texttt{#1}}

\title{Governing Equations of the \textbf{fuzzy-couscous} Solver}
\author{Clean-room compressible LES solver for blast-induced turbulence}
\date{Generated 2026-05-23}

\begin{document}
\maketitle

\begin{abstract}
This document collects every governing equation the \textbf{fuzzy-couscous}
solver discretizes: the compressible Navier--Stokes system, the inviscid and
viscous fluxes, the equations of state (ideal gas, two-$\gamma$ mixture,
Jones--Wilkins--Lee, and stiffened gas), the multifluid models (the
transported-$\gamma$ marker and the five-equation diffuse-interface model), the
LES closure
(hyperdissipation, localized artificial diffusivity, Ducros sensor), the
Besnard--Harlow--Rauenzahn (BHR) variable-density turbulence model, the SSP-RK3
time integrator, and the diagnostic identities. Each section cites the file that
implements it, so the equations stay traceable to the code.
\end{abstract}

\tableofcontents

\section{Compressible Navier--Stokes system}
\label{sec:ns}

The solver advances the conservative compressible Navier--Stokes equations for a
single fluid (or a mixture; see \S\ref{sec:eos}) on a periodic / slip-wall
Cartesian box. With conserved state
$\mathbf{U} = (\rho,\ \rho\uvec,\ \rho E)^{\mathsf T}$,
\begin{align}
\dpar{\rho}{t} + \nabla\!\cdot(\rho\uvec) &= 0, \\
\dpar{(\rho\uvec)}{t} + \nabla\!\cdot\!\left(\rho\uvec\otimes\uvec + p\,\mathbf{I}\right)
   &= \nabla\!\cdot\boldsymbol{\tau}, \\
\dpar{(\rho E)}{t} + \nabla\!\cdot\!\left[(\rho E + p)\,\uvec\right]
   &= \nabla\!\cdot(\boldsymbol{\tau}\cdot\uvec) - \nabla\!\cdot\mathbf{q},
\end{align}
where $\rho$ is density, $\uvec=(u,v,w)$ velocity, $p$ pressure, and the total
specific energy is $E = e + \tfrac12\,|\uvec|^2$ with $e$ the specific internal
energy. The right-hand side is assembled as a telescoping difference of face
fluxes, $R_i \mathrel{-}= (F_{i+1/2}-F_{i-1/2})/\Delta x$, so the discrete scheme
is conservative by construction (\code{numerics/RHS.cpp}).

\section{Inviscid flux}
\label{sec:euler}

The hyperbolic (Euler) flux in direction $d\in\{x,y,z\}$ with unit normal
$\mathbf{n}_d$ and normal velocity $u_d=\uvec\cdot\mathbf{n}_d$ is
\begin{equation}
\mathbf{F}^{\text{euler}}_d(\mathbf{U},p) =
\begin{pmatrix}
\rho u_d \\
\rho u\, u_d + p\,n_{d,x} \\
\rho v\, u_d + p\,n_{d,y} \\
\rho w\, u_d + p\,n_{d,z} \\
(\rho E + p)\,u_d
\end{pmatrix}.
\end{equation}
It is evaluated by a hybrid 6th-order central / 5th-order WENO reconstruction~\cite{jiang_shu_1996,pirozzoli_2002} in
flux-difference form with Lax--Friedrichs splitting; the WENO weight is gated by
the shock sensor of \S\ref{sec:lad} (\code{physics/EulerFlux.hpp},
\code{numerics/WENO5.hpp}). The flux is EOS-agnostic: $p$ is supplied by
\S\ref{sec:eos}.

\section{Viscous fluxes}
\label{sec:viscous}

The viscous stress is the compressible Stokes tensor with an optional bulk
viscosity $\mu_b$,
\begin{equation}
\tau_{ij} = \mu\!\left(\dpar{u_i}{x_j} + \dpar{u_j}{x_i}
            - \tfrac23\,\delta_{ij}\,\nabla\!\cdot\uvec\right)
          + \mu_b\,\delta_{ij}\,\nabla\!\cdot\uvec,
\end{equation}
and the heat flux is Fourier's law
\begin{equation}
\mathbf{q} = -\kappa\,\nabla T,
\qquad \kappa = \frac{\mu\,c_p}{\mathrm{Pr}},
\end{equation}
with $T$ the temperature and $\mathrm{Pr}$ the Prandtl number. First derivatives
use composed 6th-order central stencils (NGHOST $=6$), so the viscous Laplacian
is 6th-order accurate (\code{physics/ViscousFlux.hpp},
\code{numerics/Gradients.cpp}).

\section{Equation of state}
\label{sec:eos}

The pressure closure maps $(\text{marker},\rho,e) \mapsto (p,c)$ through a single
mixture dispatcher (\code{physics/MixtureEOS.hpp}), where $c$ is the sound speed
used at the flux and CFL sites.

\subsection{Ideal gas}
\begin{equation}
p = (\gamma-1)\,\rho e, \qquad c^2 = \frac{\gamma p}{\rho}.
\end{equation}

\subsection{Two-$\gamma$ mixture (variable-density contact)}
A per-cell advected marker $G = 1/(\gamma-1)$ selects between ambient air and a
second ideal gas (products with $\gamma_p$). The effective ratio of specific
heats is recovered as $\gamma = 1 + 1/G$, giving
\begin{equation}
p = \frac{\rho e}{G}, \qquad c^2 = \frac{(1+1/G)\,p}{\rho}.
\end{equation}
This is the original variable-density contact model (\code{physics/Multifluid.hpp}).

\subsection{Jones--Wilkins--Lee (JWL) for real high explosives}
\label{sec:jwl}
Detonation products (e.g.\ TNT) obey the JWL EOS, which an ideal gas cannot
represent. With relative volume $V = \rho_0/\rho$ ($\rho_0$ the reference/solid
density) and Gr\"uneisen coefficient $\omega$,
\begin{equation}
p(\rho,e) = A\!\left(1 - \frac{\omega}{R_1 V}\right)e^{-R_1 V}
          + B\!\left(1 - \frac{\omega}{R_2 V}\right)e^{-R_2 V}
          + \omega\,\rho\,e .
\end{equation}
Writing $f(V)$ for the reference-curve (first two) terms, $p = f(V) + \omega\rho e$,
which inverts in closed form for the internal energy at fixed density,
$e = (p - f(V))/(\omega\rho)$, and yields the isentropic sound speed
$c^2 = (\partial p/\partial\rho)_s$. With $A=B=0$ the EOS reduces to an ideal gas
with $\gamma = 1+\omega$, the hook that lets the mixture dispatcher treat air and
products uniformly (\code{physics/JWL.hpp}).

The shipped TNT parameter set (LLNL / Dobratz \& Crawford~\cite{dobratz_crawford_1985}) is:
\begin{center}
\begin{tabular}{lll}
\toprule
Parameter & Symbol & Value \\
\midrule
Reference coefficient & $A$ & \SI{371.2}{GPa} \\
Reference coefficient & $B$ & \SI{3.231}{GPa} \\
Reference exponent    & $R_1$ & $4.15$ \\
Reference exponent    & $R_2$ & $0.95$ \\
Gr\"uneisen coefficient & $\omega$ & $0.30$ \\
Reference density     & $\rho_0$ & \SI{1630}{kg/m^3} \\
Detonation energy     & $E_0$ & \SI{7.0}{GPa} \\
CJ pressure           & $p_{\mathrm{CJ}}$ & \SI{21}{GPa} \\
Detonation speed      & $D$ & \SI{6930}{m/s} \\
\bottomrule
\end{tabular}
\end{center}
Runs are nondimensionalized with $\rho_{\text{ref}}=\rho_0$ and
$p_{\text{ref}}=p_{\mathrm{CJ}}$ so the $\sim$1000:1 density and $\sim$$10^5$:1
pressure contrasts (Atwood $\approx 0.999$) stay well conditioned. The CJ density
$\rho_{\mathrm{CJ}} = \rho_0/(1 - p_{\mathrm{CJ}}/(\rho_0 D^2))$ is the
Rankine--Hugoniot value, self-consistent with the JWL sound speed at CJ
($c_{\mathrm{CJ}} = D - u_{\mathrm{CJ}}$).

\subsection{Stiffened gas}
The stiffened-gas (Noble--Abel/Tammann) EOS adds a constant stiffness pressure
$p_\infty$ to the ideal-gas law so that nearly incompressible phases (water,
condensed explosives) share the same Gr\"uneisen-type closure as a gas:
\begin{equation}
p = (\gamma-1)\,\rho e - \gamma\,p_\infty, \qquad
c^2 = \frac{\gamma\,(p + p_\infty)}{\rho}.
\end{equation}
With $p_\infty=0$ this reduces to the ideal gas (the analogue of JWL's $A=B=0$
reduction), so a stiffened-gas phase doubles as an ideal-gas phase
(\code{physics/StiffenedGas.hpp}).

\subsection{Five-equation mixture (volume-fraction / Wood rule)}
\label{sec:eos5eq}
For the five-equation model (\S\ref{sec:fiveeq})~\cite{allaire_2002,murrone_guillard_2005} two phases coexist in each cell
with volume fraction $\alpha_1$ ($\alpha_2 = 1-\alpha_1$) and partial densities
$Z_k = \alpha_k\rho_k$. Each phase obeys a Gr\"uneisen split of its
internal-energy density, $\rho_k e_k = (p + \Phi_k)/(\Gamma_k-1)$, where
\begin{equation}
\Gamma_k - 1 =
\begin{cases} \gamma_k - 1 & \text{stiffened gas} \\ \omega_k & \text{JWL} \end{cases},
\qquad
\Phi_k =
\begin{cases} \gamma_k\,p_{\infty,k} & \text{stiffened gas} \\ -f_k(V_k) & \text{JWL (ref.\ curve, \S\ref{sec:jwl})} \end{cases}.
\end{equation}
Summing $\rho e = \sum_k \alpha_k \rho_k e_k$ over the phases gives the
volume-fraction-averaged (Wood) pressure
\begin{equation}
p = \frac{\rho e - \Pi}{S}, \qquad
S = \sum_k \frac{\alpha_k}{\Gamma_k-1}, \qquad
\Pi = \sum_k \frac{\alpha_k\,\Phi_k}{\Gamma_k-1},
\end{equation}
with $\rho e$ the mixture internal-energy density and $\rho = Z_1+Z_2$. The
frozen mixture sound speed is the mass-weighted phase average
\begin{equation}
c^2 = \sum_k Y_k\,c_k^2, \qquad Y_k = Z_k/\rho,
\end{equation}
where $c_k^2$ is the per-phase stiffened-gas/JWL sound speed at the common
pressure $p$. For stiffened-gas phases $S$ and $\Pi$ are linear in $\alpha_1$ ---
the property that lets an isolated interface stay pressure-oscillation-free
(\code{physics/MixtureEOS.hpp}, \code{p\_c\_5eq}).

\section{Multifluid models}
\label{sec:multifluid}

\subsection{Model taxonomy}
The solver carries two multifluid models, both single-velocity and
single-pressure, related by the diffuse-interface hierarchy (Baer--Nunziato's
seven-equation non-equilibrium system~\cite{baer_nunziato_1986} being the full parent):
\begin{itemize}
\item \textbf{Transported-$\gamma$ four-equation model}~\cite{shyue_1998,abgrall_karni_2001,karni_1994,abgrall_1996}:
  the three conservative hydrodynamic equations of
  \S\ref{sec:ns} plus \emph{one} advected EOS marker (below). The equation of
  state is selected per cell by the marker; there are no per-phase masses. This
  is distinct from the \emph{mass-fraction} four-equation model
  (two partial masses closed by temperature equilibrium~\cite{murrone_guillard_2005}), which this code does
  not implement.
\item \textbf{Five-equation diffuse-interface model}~\cite{allaire_2002,kapila_2001}:
  velocity- and pressure-equilibrium two-phase flow with two partial-mass
  equations, mixture momentum and energy, and a volume-fraction transport
  (\S\ref{sec:fiveeq}).
\end{itemize}
The $4\!\to\!5$ transition splits the single mixture mass into the two partial
masses $\alpha_1\rho_1,\alpha_2\rho_2$ and replaces the transported $\gamma$
marker by the volume fraction $\alpha_1$.

\subsection{Transported-$\gamma$ marker transport}
The marker (the two-$\gamma$ field $G$, or the products mass fraction
$\phi\in[0,1]$ for JWL) is advected with the flow,
\begin{equation}
\dpar{(\rho\,\chi)}{t} + \nabla\!\cdot(\rho\,\chi\,\uvec) = 0,
\qquad \chi \in \{G,\ \phi\},
\end{equation}
equivalently $\partial\chi/\partial t + \uvec\cdot\nabla\chi = 0$ by continuity.
The contact is handled either by a gated double-flux (non-oscillatory) or an
opt-in fully conservative flux (\code{conservative = true}); the latter conserves
total energy at the cost of small contact oscillations
(\code{physics/Multifluid.cpp}). The double-flux method --- freezing the per-cell
EOS across the flux stencil so a single-material conservative scheme handles the
contact without spurious pressure oscillation --- is that of
Abgrall \& Karni~\cite{abgrall_karni_2001} (building on
Karni~\cite{karni_1994} and Abgrall~\cite{abgrall_1996}); ``gated'' here denotes
that it is applied only within a few cells of a contact (bulk cells take the
cheap shared flux), an implementation optimization rather than a separately
named scheme.

\subsection{Five-equation diffuse-interface model}
\label{sec:fiveeq}
The true five-equation system~\cite{allaire_2002,kapila_2001} advances the two
partial masses, the mixture momentum and energy, and the volume fraction:
\begin{align}
\dpar{(\alpha_1\rho_1)}{t} + \nabla\!\cdot(\alpha_1\rho_1\,\uvec) &= 0, \\
\dpar{(\alpha_2\rho_2)}{t} + \nabla\!\cdot(\alpha_2\rho_2\,\uvec) &= 0, \\
\dpar{(\rho\uvec)}{t} + \nabla\!\cdot\!\left(\rho\uvec\otimes\uvec + p\,\mathbf{I}\right)
   &= \nabla\!\cdot\boldsymbol{\tau}, \\
\dpar{(\rho E)}{t} + \nabla\!\cdot\!\left[(\rho E + p)\,\uvec\right]
   &= \nabla\!\cdot(\boldsymbol{\tau}\cdot\uvec) - \nabla\!\cdot\mathbf{q}, \\
\dpar{\alpha_1}{t} + \uvec\cdot\nabla\alpha_1 &= 0,
\end{align}
with mixture density $\rho = \alpha_1\rho_1 + \alpha_2\rho_2$ and pressure $p$
from the Wood-rule law (\S\ref{sec:eos5eq}). The non-conservative volume-fraction
equation is discretized in quasi-conservative form,
$\partial\alpha_1/\partial t + \nabla\!\cdot(\alpha_1\uvec) = \alpha_1\,\nabla\!\cdot\uvec$,
using the same face velocities as the conservative fluxes. The two partial-mass
fluxes reuse the same reconstruction as the mixture-mass flux, so $Z_1+Z_2\equiv
\rho$ to round-off, and the interface stays pressure/velocity-equilibrium
(no spurious oscillation at a contact). The state is advanced in lockstep with
the conserved fields through every SSP-RK3 stage
(\code{numerics/RHS.cpp}, \code{add\_inviscid\_divergence\_5eq};
\code{numerics/RK3.cpp}, \code{step\_5eq}/\code{step\_mpi\_5eq}).

\section{LES closure and shock capturing}
\label{sec:lad}

\subsection{Spectral hyperdissipation (LES sink)}
The sub-grid sink is a hyper-viscous term applied to every conserved variable,
\begin{equation}
\left.\dpar{\mathbf{U}}{t}\right|_{\text{LES}} = -\,\nu_h\,\nabla^4 \mathbf{U},
\end{equation}
whose discrete operator matches the analytic $k^4$ damping; under SSP-RK3 it
preserves the $e^{-\nu_h k^4 t}$ eigenvalue decay (\code{numerics/HyperdissipationSpectral.cpp}).
The empirical stability invariant is fixed $\nu_{\text{eff}}=\nu_h k_{\max}^2$;
between resolutions $\nu_h$ scales as $\Delta x^2$ (see the README TGV note).

\subsection{Ducros shock sensor}
The hybrid central/WENO blend~\cite{pirozzoli_2002} and the artificial diffusivity are gated by the
Ducros sensor~\cite{ducros_1999}
\begin{equation}
\theta = \frac{(\nabla\!\cdot\uvec)^2}{(\nabla\!\cdot\uvec)^2 + |\boldsymbol{\omega}|^2 + \varepsilon},
\end{equation}
with vorticity $\boldsymbol{\omega}=\nabla\times\uvec$ and a small $\varepsilon$;
$\theta\to1$ in compression (shocks), $\theta\to0$ in smooth/vortical regions. A
pressure-jump $\tanh$ ramp and a half-width-2 dilation along the face direction
sharpen the gate (\code{numerics/Ducros.cpp}).

\subsection{Localized artificial diffusivity (LAD)}
For strong blasts the Cook/Kawai--Lele localized artificial diffusivity~\cite{cook_2007,kawai_lele_2008} adds
grid-dependent, sensor-localized dissipation to shocks and contacts,
\begin{equation}
\mu^\ast = C_\mu\,\rho\,|\nabla^r(\nabla\!\cdot\uvec)|\,\Delta^{r+2},
\qquad
\beta^\ast = C_\beta\,\rho\,|\nabla^r(\nabla\!\cdot\uvec)|\,\Delta^{r+2},
\end{equation}
with an analogous artificial mass/contact diffusivity $D^\ast$ for the
density-contact piece; the canonical order is $r=4$ (\code{numerics/ArtificialDiffusivity.cpp}).

\section{BHR variable-density turbulence model}
\label{sec:bhr}

The Besnard--Harlow--Rauenzahn~\cite{besnard_1992} $k$--$\varepsilon$--$a$--$b$ model runs as an
operator-split sub-step alongside the LES hydro. The six primitives are the
turbulent kinetic energy $k$, its dissipation $\varepsilon$, the turbulent mass
flux $a_i$, and the density--specific-volume covariance
$b = -\,\overline{\rho'\,(1/\rho)'}$. The dominant TKE source is the
variable-density (baroclinic) production
\begin{equation}
P_{\mathrm{VD}} = -\,a_i\,\dpar{p}{x_i},
\end{equation}
validated as dominant in the LES $a/b$ budget. The turbulent viscosity is
$\nu_t = C_\mu\,k^2/\varepsilon$, and the modeled transport equations carry
gradient-diffusion terms with Schmidt numbers
$(\sigma_k,\sigma_\varepsilon,\sigma_a,\sigma_b)$
(\code{turbulence/BHR.hpp}, \code{turbulence/BHR.cpp}).

When two-way feedback is enabled, the model returns a Reynolds-stress divergence
and dissipative heating to $\mathbf{U}$, blended by the resolution-adequacy sensor
\begin{equation}
f_k = \frac{k_{\mathrm{sgs}}}{k_{\mathrm{res}} + k_{\mathrm{sgs}}},
\qquad
k_{\mathrm{res}} = \tfrac12\,\big|\uvec - \widetilde{\uvec}\big|^2,
\end{equation}
where $\widetilde{\uvec}$ is a test-filtered velocity (resolved energy in the
$[\Delta x, 2\Delta x]$ octave); $f_k\to1$ under-resolved (RANS), $f_k\to0$ when
the LES resolves the energy.

\section{Time integration}
\label{sec:rk3}

Time advance uses the three-stage strong-stability-preserving Runge--Kutta
(SSP-RK3, Gottlieb--Shu~\cite{gottlieb_shu_1998}) scheme on $\mathbf{U}^n \to \mathbf{U}^{n+1}$ with
right-hand side $\mathcal{L}(\mathbf{U})$:
\begin{align}
\mathbf{U}^{(1)} &= \mathbf{U}^n + \Delta t\,\mathcal{L}(\mathbf{U}^n),\\
\mathbf{U}^{(2)} &= \tfrac34\,\mathbf{U}^n + \tfrac14\,\mathbf{U}^{(1)}
                   + \tfrac14\,\Delta t\,\mathcal{L}(\mathbf{U}^{(1)}),\\
\mathbf{U}^{n+1} &= \tfrac13\,\mathbf{U}^n + \tfrac23\,\mathbf{U}^{(2)}
                   + \tfrac23\,\Delta t\,\mathcal{L}(\mathbf{U}^{(2)}).
\end{align}
The time step obeys the minimum of a hyperbolic CFL ($\propto \Delta x/(|\uvec|+c)$)
and a viscous limit; in MPI runs both are reduced with \code{MPI\_Allreduce(MIN)}
(\code{numerics/RK3.cpp}, \code{numerics/RHS.cpp}).

\section{Diagnostic identities}
\label{sec:diag}

\subsection{Helmholtz decomposition}
The velocity is split into solenoidal and dilatational parts,
$\uvec = \uvec_s + \uvec_d$ with $\nabla\!\cdot\uvec_s = 0$ and
$\nabla\times\uvec_d = \mathbf{0}$, giving the kinetic-energy split
$K = K_{\mathrm{sol}} + K_{\mathrm{dil}}$ and the per-shell spectra
$E(k) = E_{\mathrm{sol}}(k) + E_{\mathrm{dil}}(k)$ (\code{diagnostics/Spectra.cpp}).

\subsection{Dissipation budget}
The resolved dissipation is split into solenoidal and dilatational contributions,
\begin{equation}
\varepsilon_{\text{total}} = \frac{\mu}{\rho}\,\big\langle \tau_{ij}\,\partial_j u_i\big\rangle,
\qquad
\varepsilon_{\mathrm{sol}} = \nu\,\langle|\boldsymbol{\omega}|^2\rangle,
\qquad
\varepsilon_{\mathrm{dil}} = \tfrac43\,\nu\,\langle(\nabla\!\cdot\uvec)^2\rangle,
\end{equation}
together with the turbulent Mach number $M_t = u_{\text{rms}}/\langle c\rangle$
(\code{diagnostics/Statistics.cpp}).

\subsection{Conservation monitors}
For a closed chamber the totals $E/E_0$ and $M/M_0$ are conserved to round-off by
the discrete divergence theorem (telescoping face fluxes plus $\uvec\cdot\mathbf{n}=0$
at slip walls); the internal energy $e_{\text{int}} = \sum(\rho E - \tfrac12\rho|\uvec|^2)$
is accumulated directly so $KE + e_{\text{int}} = e_{\text{total}}$ is a genuine
consistency check rather than an identity (\code{diagnostics/Statistics.cpp};
see the README for the conservation-vs-truncation discussion).

\begin{thebibliography}{99}

\bibitem{baer_nunziato_1986}
M.~R. Baer and J.~W. Nunziato,
``A two-phase mixture theory for the deflagration-to-detonation transition (DDT)
in reactive granular materials,''
\textit{Int.\ J.\ Multiphase Flow} \textbf{12}(6):861--889, 1986.

\bibitem{karni_1994}
S.~Karni,
``Multicomponent flow calculations by a consistent primitive algorithm,''
\textit{J.\ Comput.\ Phys.} \textbf{112}(1):31--43, 1994.

\bibitem{abgrall_1996}
R.~Abgrall,
``How to prevent pressure oscillations in multicomponent flow calculations:
a quasi conservative approach,''
\textit{J.\ Comput.\ Phys.} \textbf{125}(1):150--160, 1996.

\bibitem{jiang_shu_1996}
G.-S. Jiang and C.-W. Shu,
``Efficient implementation of weighted ENO schemes,''
\textit{J.\ Comput.\ Phys.} \textbf{126}(1):202--228, 1996.

\bibitem{shyue_1998}
K.-M. Shyue,
``An efficient shock-capturing algorithm for compressible multicomponent
problems,''
\textit{J.\ Comput.\ Phys.} \textbf{142}(1):208--242, 1998.

\bibitem{gottlieb_shu_1998}
S.~Gottlieb and C.-W. Shu,
``Total variation diminishing Runge--Kutta schemes,''
\textit{Math.\ Comp.} \textbf{67}(221):73--85, 1998.

\bibitem{ducros_1999}
F.~Ducros, V.~Ferrand, F.~Nicoud, C.~Weber, D.~Darracq, C.~Gacherieu, and
T.~Poinsot,
``Large-eddy simulation of the shock/turbulence interaction,''
\textit{J.\ Comput.\ Phys.} \textbf{152}(2):517--549, 1999.

\bibitem{kapila_2001}
A.~K. Kapila, R.~Menikoff, J.~B. Bdzil, S.~F. Son, and D.~S. Stewart,
``Two-phase modeling of deflagration-to-detonation transition in granular
materials: reduced equations,''
\textit{Phys.\ Fluids} \textbf{13}(10):3002--3024, 2001.

\bibitem{abgrall_karni_2001}
R.~Abgrall and S.~Karni,
``Computations of compressible multifluids,''
\textit{J.\ Comput.\ Phys.} \textbf{169}(2):594--623, 2001.

\bibitem{pirozzoli_2002}
S.~Pirozzoli,
``Conservative hybrid compact-WENO schemes for shock-turbulence interaction,''
\textit{J.\ Comput.\ Phys.} \textbf{178}(1):81--117, 2002.

\bibitem{allaire_2002}
G.~Allaire, S.~Clerc, and S.~Kokh,
``A five-equation model for the simulation of interfaces between compressible
fluids,''
\textit{J.\ Comput.\ Phys.} \textbf{181}(2):577--616, 2002.

\bibitem{murrone_guillard_2005}
A.~Murrone and H.~Guillard,
``A five equation reduced model for compressible two phase flow problems,''
\textit{J.\ Comput.\ Phys.} \textbf{202}(2):664--698, 2005.

\bibitem{cook_2007}
A.~W. Cook,
``Artificial fluid properties for large-eddy simulation of compressible
turbulent mixing,''
\textit{Phys.\ Fluids} \textbf{19}(5):055103, 2007.

\bibitem{kawai_lele_2008}
S.~Kawai and S.~K. Lele,
``Localized artificial diffusivity scheme for discontinuity capturing on
curvilinear meshes,''
\textit{J.\ Comput.\ Phys.} \textbf{227}(22):9498--9526, 2008.

\bibitem{besnard_1992}
D.~Besnard, F.~H. Harlow, R.~M. Rauenzahn, and C.~Spangler,
``Turbulence transport equations for variable-density turbulence and their
relationship to two-field models,''
Tech.\ Rep.\ LA-12303-MS, Los Alamos National Laboratory, 1992.

\bibitem{dobratz_crawford_1985}
B.~M. Dobratz and P.~C. Crawford,
``LLNL Explosives Handbook: Properties of Chemical Explosives and Explosive
Simulants,''
Tech.\ Rep.\ UCRL-52997, Lawrence Livermore National Laboratory, 1985.

\end{thebibliography}

\end{document}
